% \iffalse meta-comment
%
% hit-outline.dtx 是照轮廓画的字标与美术字
%
% \fi
%
% \section{照轮廓画的字标与美术字}
%
% \changes{v3.2a}{2026/09/27}{\hologo{pdfLaTeX} 下用 \cs{pdfliteral} 画（xdvipdfmx 的 special 在 pdfTeX 下被忽略）；逐字查轮廓
%  按字取，\hologo{pdfLaTeX} 下一个汉字是几个记号}
% \changes{v3.2a}{2026/09/26}{报告首页的校名字标、本科报告落款的隶书、深圳本科报告首页的
%   华文新魏照 iota-hit 的轮廓直接画：字标不再随包带 PDF；两种字体机器上没有时照轮廓画，
%   不再贴图或退成宋体}
%
% 封面上有三样东西正文字体排不出来：报告首页的校名字标（学校给的书法字，本来就是图）、
% 本科报告落款“哈尔滨工业大学教务处制”（隶书）、深圳本科报告首页的校名与文种（华文新魏）。
% 后两种字体只有 Windows 与装了 Office 的机器上有。先前字标与落款各贴一份 PDF，
% 新魏没有就退成宋体。
%
% 现在三样都照轮廓在页面上画。轮廓取自 iota-hit：字标是它从学校的 SVG 抽的，每个字一段路径；
% 隶书与新魏是它从字体里抽的那十几个字。\file{scripts/iota-outlines.py} 把这些数据换成 PDF 的
% 路径算子，写进 \file{src/assets/hit-outline-data.dtx}，由它生成 \file{hithesis-outline.def}。
% 那边的数据改了，再跑一遍脚本。
%
% 隶书与新魏在机器上有的时候照旧排字体，走哪条路由 \option{cover-script} 定，见
% \file{src/engine/hit-fonts.dtx}。
%
% \emph{一个字一只盒。} 盒宽是字体的步进，高与深是字体的上升部与下降部（iota-hit 数据里的
% \texttt{asc}、\texttt{desc}，隶书每 em 256、上升 220、下降 36，与中易宋体一样）。基线上放一段
% PDF literal，缩放到字号再填。加粗照 \pkg{luaotfload} 合成粗体的描边：线宽是字号的 1/35，
% 字体那条路量出来 18\,bp 落款的线宽是 0.515、36\,bp 新魏 1.030，正是这个数。
%
% 画出来的字没有文字层，复制与搜索都拿不到。每只盒套一段带 \texttt{ActualText} 的标记内容，
% 读屏与认它的阅读器能读出原字。
%
% \emph{字标整只画成一只盒}，用 \cs{docxlike\_picture:n} 标成嵌入型图片，所在那一段照装图段
% 定行高，与先前 \cs{includegraphics} 插的 PDF 一样。
%
% \emph{不慢。} 画 200 次字标比 \cs{includegraphics} 插 200 次 PDF 还快（0.32 秒对 0.44 秒，
% 空文档 0.26 秒）。数据一百多 KB，第一次用到时才读，用不着的文档一次都不读。
%
%    \begin{macrocode}
%<*shared-outline>
\bool_new:N \g__hit_outline_loaded_bool
\bool_new:N \l__hit_outline_bold_bool
\bool_new:N \l__hit_outline_cover_bool
\dim_new:N \l__hit_outline_size_dim
\dim_new:N \l__hit_outline_extra_dim
\dim_new:N \l__hit_outline_height_dim
\tl_new:N \l__hit_outline_text_tl
\tl_new:N \l__hit_outline_key_tl
\tl_new:N \l__hit_outline_glyph_tl
\tl_new:N \l__hit_outline_scale_tl
\tl_new:N \l__hit_outline_stroke_tl
\tl_new:N \l__hit_outline_upm_tl
\tl_new:N \l__hit_outline_asc_tl
\tl_new:N \l__hit_outline_desc_tl
\box_new:N \l__hit_outline_box
%    \end{macrocode}
%
% 读数据。数据是按 \pkg{expl3} 语法写的（空格写成 \texttt{\textasciitilde}），在组里打开
% 这套语法读，读完还回去。文件找不到警告一次，字标空着，隶书与新魏排字体或宋体。
%    \begin{macrocode}
\cs_new_protected:Npn \__hit_outline_load:
  {
    \bool_if:NF \g__hit_outline_loaded_bool
      {
        \bool_gset_true:N \g__hit_outline_loaded_bool
        \file_if_exist:nTF { hithesis-outline.def }
          {
            \group_begin:
              \ExplSyntaxOn
              \file_input:n { hithesis-outline.def }
            \group_end:
          }
          { \msg_warning:nn { hithesis } { outline-missing } }
      }
  }
%    \end{macrocode}
%
% PDF literal。三个引擎都把原点挪到当前点，单位是 bp。
%    \begin{macrocode}
\sys_if_engine_luatex:TF
  {
    \cs_new_protected:Npn \__hit_outline_literal:e #1
      { \tex_pdfextension:D literal {#1} }
  }
  {
    \sys_if_engine_pdftex:TF
      {
        \cs_new_protected:Npn \__hit_outline_literal:e #1
          { \tex_pdfliteral:D {#1} }
      }
      {
        \cs_new_protected:Npn \__hit_outline_literal:e #1
          { \tex_special:D { pdf:literal~#1 } }
      }
  }
%    \end{macrocode}
%
% \begin{macro}{\hit@outline@if@cover:nnTF}
% \marg{字体}\marg{文本}：文本里的字是不是都有轮廓。字体写 \texttt{lishu} 或 \texttt{xinwei}。
% 文本全展开再逐字查；展开后还剩命令的，查不到，算没有。
%    \begin{macrocode}
\prg_new_protected_conditional:Npnn \hit@outline@if@cover:nn #1#2 { T , F , TF }
  {
    \__hit_outline_load:
    \tl_set:Ne \l__hit_outline_text_tl {#2}
    \hit@engine@group@chars:N \l__hit_outline_text_tl
    \prop_if_exist:cTF { c__hit_outline_#1_prop }
      {
        \bool_set_true:N \l__hit_outline_cover_bool
        \tl_map_inline:Nn \l__hit_outline_text_tl
          {
            \hit@engine@char@key:nN {##1} \l__hit_outline_key_tl
            \prop_if_in:cVF { c__hit_outline_#1_prop } \l__hit_outline_key_tl
              {
                \bool_set_false:N \l__hit_outline_cover_bool
                \tl_map_break:
              }
          }
        \bool_if:NTF \l__hit_outline_cover_bool
          { \prg_return_true: } { \prg_return_false: }
      }
      { \prg_return_false: }
  }
\prg_generate_conditional_variant:Nnn \hit@outline@if@cover:nn { nV } { T , F , TF }
%    \end{macrocode}
% \end{macro}
%
% \begin{macro}{\hit@outline:nn}
% \marg{字体}\marg{文本}：照轮廓一个字一只盒排出来，字号与粗细取当前字体（\cs{f@size}、
% \cs{f@series}）。没有轮廓的字原样排。\cs{l\_\_hit\_outline\_extra\_dim} 是每个字后面补的
% 宽，默认零，要补的地方自己设，见 \cs{hit@cover@bottommark}。
%    \begin{macrocode}
\cs_new_protected:Npn \hit@outline:nn #1#2
  {
    \__hit_outline_load:
    \prop_if_exist:cTF { c__hit_outline_#1_prop }
      {
        \dim_set:Nn \l__hit_outline_size_dim { \f@size pt }
        \str_if_in:NnTF \f@series { b }
          { \bool_set_true:N \l__hit_outline_bold_bool }
          { \bool_set_false:N \l__hit_outline_bold_bool }
        \exp_last_unbraced:Nv \__hit_outline_metrics:nnn { c__hit_outline_#1_metrics_tl }
        \tl_set:Ne \l__hit_outline_text_tl {#2}
        \hit@engine@group@chars:N \l__hit_outline_text_tl
        \tl_map_inline:Nn \l__hit_outline_text_tl
          {
            \hit@engine@char@key:nN {##1} \l__hit_outline_key_tl
            \prop_get:cVNTF { c__hit_outline_#1_prop } \l__hit_outline_key_tl \l__hit_outline_glyph_tl
              {
                \exp_after:wN \__hit_outline_glyph:nnn \l__hit_outline_glyph_tl
                \dim_compare:nNnF \l__hit_outline_extra_dim = \c_zero_dim
                  { \tex_kern:D \l__hit_outline_extra_dim }
              }
              {##1}
          }
      }
      {#2}
  }
\cs_new_protected:Npn \__hit_outline_metrics:nnn #1#2#3
  {
    \tl_set:Nn \l__hit_outline_upm_tl {#1}
    \tl_set:Nn \l__hit_outline_asc_tl {#2}
    \tl_set:Nn \l__hit_outline_desc_tl {#3}
    \tl_set:Ne \l__hit_outline_scale_tl
      {
        \fp_eval:n
          { round ( \dim_to_decimal_in_bp:n { \l__hit_outline_size_dim } / #1 , 6 ) }
      }
    \tl_set:Ne \l__hit_outline_stroke_tl { \fp_eval:n { round ( #1 / 35 , 3 ) } }
  }
%    \end{macrocode}
% \end{macro}
%
% 一个字：\marg{步进}\marg{UTF-16}\marg{路径}，路径在字体单位里，基线在零、$y$ 向上，
% 与 PDF 同向，只缩放不翻转。粗体填了再描（\texttt{B}），线宽也在字体单位里，跟着缩放。
%    \begin{macrocode}
\cs_new_protected:Npn \__hit_outline_glyph:nnn #1#2#3
  {
    \hbox_set:Nn \l__hit_outline_box
      {
        \__hit_outline_literal:e
          {
            /Span << /ActualText <FEFF #2> >> BDC ~
            q ~ \l__hit_outline_scale_tl \c_space_tl 0 ~ 0 ~
            \l__hit_outline_scale_tl \c_space_tl 0 ~ 0 ~ cm ~
            \bool_if:NTF \l__hit_outline_bold_bool
              { \l__hit_outline_stroke_tl \c_space_tl w ~ 0 ~ j ~ 10 ~ M ~ #3 ~ B ~ }
              { #3 ~ f ~ }
            Q ~ EMC
          }
      }
    \box_set_wd:Nn \l__hit_outline_box
      { \l__hit_outline_size_dim * #1 / \l__hit_outline_upm_tl }
    \box_set_ht:Nn \l__hit_outline_box
      { \l__hit_outline_size_dim * \l__hit_outline_asc_tl / \l__hit_outline_upm_tl }
    \box_set_dp:Nn \l__hit_outline_box
      {
        \l__hit_outline_size_dim * \int_eval:n { 0 - \l__hit_outline_desc_tl }
          / \l__hit_outline_upm_tl
      }
    \box_use_drop:N \l__hit_outline_box
  }
%    \end{macrocode}
%
% \begin{macro}{\hit@wordmark:nn}
% \marg{宽}\marg{高}：校名字标，宽高是原件里图框的大小。原件嵌的是一张 466 $\times$ 87 的
% JPEG，框是 167.8 $\times$ 31.25\,磅（VML 的 \texttt{style}）；iota-hit 的轮廓按学校 SVG 的
% 516 比 100 画，同宽时高 32.52。不压扁，照 iota-hit 从框顶往下画，墨迹只到视框高的 91\%，
% 正好落在框里。与 Windows Word 的导出件比，墨迹左边差 0.24、右边 0.72、上沿 0.6\,磅，下沿
% 重合；先前那份 PDF 的墨迹左右各宽出将近 5\,磅。
%    \begin{macrocode}
\cs_new_protected:Npn \hit@wordmark:nn #1#2
  {
    \__hit_outline_load:
    \tl_if_exist:NT \c__hit_outline_wordmark_tl
      {
        \dim_set:Nn \l__hit_outline_size_dim {#1}
        \dim_set:Nn \l__hit_outline_height_dim {#2}
        \exp_last_unbraced:NV \__hit_outline_wordmark:nn \c__hit_outline_wordmark_box_tl
      }
  }
\cs_new_protected:Npn \__hit_outline_wordmark:nn #1#2
  {
    \tl_set:Ne \l__hit_outline_scale_tl
      {
        \fp_eval:n
          { round ( \dim_to_decimal_in_bp:n { \l__hit_outline_size_dim } / #1 , 6 ) }
      }
    \hbox_set:Nn \l__hit_outline_box
      {
        \__hit_outline_literal:e
          {
            /Span << /ActualText <FEFF \c__hit_outline_wordmark_text_tl > >> BDC ~
            q ~ \l__hit_outline_scale_tl \c_space_tl 0 ~ 0 ~
            \l__hit_outline_scale_tl \c_space_tl 0 ~
            \fp_eval:n
              {
                round ( \dim_to_decimal_in_bp:n { \l__hit_outline_height_dim }
                  - #2 * \l__hit_outline_scale_tl , 3 )
              }
            \c_space_tl cm ~
            \c__hit_outline_wordmark_tl \c_space_tl f ~ Q ~ EMC
          }
      }
    \box_set_wd:Nn \l__hit_outline_box { \l__hit_outline_size_dim }
    \box_set_ht:Nn \l__hit_outline_box { \l__hit_outline_height_dim }
    \box_set_dp:Nn \l__hit_outline_box { \c_zero_dim }
    \docxlike_picture:n { \box_use_drop:N \l__hit_outline_box }
  }
%    \end{macrocode}
% \end{macro}
%
%    \begin{macrocode}
%</shared-outline>
%    \end{macrocode}
%
% \Finale
